\documentclass[10pt,letterpaper]{article}
\usepackage[letterpaper,margin=0.72in,headheight=20pt,footskip=28pt]{geometry}
\usepackage[T1]{fontenc}\usepackage{lmodern}\usepackage[table]{xcolor}\usepackage{array,longtable,booktabs,tabularx}\usepackage{enumitem}\usepackage[most]{tcolorbox}\usepackage{fancyhdr,lastpage}\usepackage{microtype}\usepackage{amsmath,amssymb}\usepackage[hidelinks,bookmarks=false]{hyperref}
\definecolor{navy}{HTML}{12304A}\definecolor{teal}{HTML}{137C8B}\definecolor{cuebg}{HTML}{EAF2F5}\definecolor{notebg}{HTML}{F8FAFB}\definecolor{line}{HTML}{B8C8D1}\definecolor{gold}{HTML}{C9922E}
\setlength{\parindent}{0pt}\setlength{\parskip}{4pt}\setlist[itemize]{leftmargin=1.25em,itemsep=1pt,topsep=2pt}\setlist[enumerate]{leftmargin=1.5em,itemsep=1pt,topsep=2pt}\renewcommand{\arraystretch}{1.18}
\pagestyle{fancy}\fancyhf{}\lhead{\small\color{navy}\textbf{CISSP CBK Cornell Notes}}\rhead{\small\color{teal}Domain 8: Software Development Security}\cfoot{\small\color{navy}Page \thepage\ of \pageref{LastPage}}\renewcommand{\headrulewidth}{0.4pt}\renewcommand{\headrule}{\hbox to\headwidth{\color{teal}\leaders\hrule height \headrulewidth\hfill}}
\newenvironment{cornell}{\rowcolors{2}{white}{notebg}\begin{longtable}{@{}>{\raggedright\arraybackslash\columncolor{cuebg}\color{navy}\bfseries}p{0.285\textwidth}>{\raggedright\arraybackslash}p{0.65\textwidth}@{}}\rowcolor{navy}\color{white}\textbf{Cue / Active Recall} & \color{white}\textbf{Detailed Notes}\\ \hline\endfirsthead\rowcolor{navy}\color{white}\textbf{Cue / Active Recall (cont.)} & \color{white}\textbf{Detailed Notes}\\ \hline\endhead}{\end{longtable}}
\newcommand{\noterow}[2]{#1 & #2\\\arrayrulecolor{line}\hline}\newcommand{\source}[1]{\par{\footnotesize\color{teal}\textit{Source: CBK, #1}}\par}
\newtcolorbox{summarybox}[1]{breakable,colback=cuebg,colframe=teal,boxrule=0.7pt,arc=2mm,title=\textbf{#1},fonttitle=\color{white},colbacktitle=teal}\newtcolorbox{exambox}[1]{breakable,colback=white,colframe=gold,boxrule=0.7pt,arc=2mm,title=\textbf{#1},fonttitle=\color{white},colbacktitle=gold}
\begin{document}\hypersetup{pageanchor=false}
\begin{titlepage}\pagecolor{navy}\color{white}\vspace*{1.0in}{\Large THE OFFICIAL (ISC)\textsuperscript{2} CISSP CBK REFERENCE\par}\vspace{0.25in}{\large Sixth Edition \textbullet\ 2021\par}\vfill{\Huge\bfseries Cornell Notes\par}\vspace{0.25in}{\LARGE Domain 8\par}\vspace{0.08in}{\LARGE\bfseries Software Development Security\par}\vfill{\large Arthur Deane and Aaron Kraus\par}\vspace{0.12in}{\normalsize Study and review companion\par}\vspace{0.35in}{\small Primary source: \textit{The Official (ISC)\textsuperscript{2} CISSP CBK Reference}, Sixth Edition, pp. 549--624.\par}\end{titlepage}
\hypersetup{pageanchor=true}\nopagecolor\color{black}

\section*{\color{navy}Domain Overview}
Software security is lifecycle security. The chapter moves from development methodologies and maturity, through the ecosystem that creates and deploys code, to effectiveness assessment, acquired-software risk, secure coding, API protection, and software-defined controls. The recurring CISSP idea is to integrate security early, measure it continuously, and retain governance even when development, infrastructure, or software comes from another party.

\begin{summarybox}{Review Objectives}
Integrate security through the SDLC/SLC; distinguish Waterfall, Agile-family methods, DevOps/DevSecOps, Spiral, Prototype, and Cleanroom; use maturity models; secure development ecosystems and CI/CD; compare SAST/DAST/IAST/RASP; audit and analyze software risk; assess COTS/OSS/third-party/cloud software; apply CWE/CVSS/OWASP concepts, API safeguards, secure coding practices, and software-defined security.
\end{summarybox}

\section{\color{navy}Security in the Software Development Life Cycle}
\source{pp. 550--560}
\begin{cornell}
\noterow{SDLC vs. SLC}{The \textbf{SDLC} organizes activities used to create and evolve software. The broader \textbf{system lifecycle (SLC)} continues through operations/maintenance and decommissioning/disposal. Security requirements, resources, verification, vulnerability management, and eventual sanitization must follow the applicable lifecycle phases.}
\noterow{Common lifecycle flow}{The chapter groups lifecycle work as initiation; development; deployment/delivery; operations/maintenance; and disposal. Security is an engineering concern across these activities, not a standalone gate. NIST SP 800-160 aligns systems-security engineering with the broader lifecycle approach in ISO/IEC/IEEE 15288.}
\noterow{Waterfall}{Sequential requirements $\rightarrow$ design $\rightarrow$ implementation $\rightarrow$ testing $\rightarrow$ deployment $\rightarrow$ maintenance. A phase is completed before advancing. Rigor can strengthen requirements and testing, but late change is difficult; the model fits stable, knowable requirements better than rapidly changing ones.}
\noterow{Agile}{Agile is a set of guiding principles, not one methodology. It emphasizes customer value, early/continuous delivery, cross-functional communication, short iterations, and acceptance of changing requirements. Security stakeholders must participate so each iteration preserves security requirements and tests new functionality.}
\noterow{Agile testing}{Testing moves close to development and is repeated on smaller changes. Automation supports frequent functional, code, regression, and vulnerability tests and is an enabler of continuous integration/continuous delivery. The aim is rapid feedback before defects reach production.}
\noterow{Extreme Programming (XP)}{An Agile method built around high-value user stories and close stakeholder participation. \textbf{Pair programming} adds continuous peer input; \textbf{refactoring} removes/reworks obsolete or unnecessary code (cruft) while preserving required behavior.}
\noterow{Test-Driven Development (TDD)}{Write a test first, run it, then write/refactor code until it passes. Tests make desired behavior explicit and support controlled modification of existing systems, including code-level security fixes.}
\noterow{Scrum roles}{\textbf{Scrum master}: facilitates and removes impediments. \textbf{Product owner}: represents product/customer value and prioritizes requirements. \textbf{Development team}: cross-functional group that delivers the increment, including delegated testing/security work.}
\noterow{Scrum artifacts / cadence}{The product \textbf{backlog} contains prioritized work; a \textbf{sprint} is a time-boxed set of committed work; sprint planning/review/retrospective surround delivery; the \textbf{daily Scrum} (standup) synchronizes progress, impediments, and near-term work. Changes normally wait for a future sprint once work is committed.}
\noterow{DevOps}{Integrates development and operations into a cross-functional delivery capability. Automation, frequent testing, pipelines, integration, deployment, and QA reduce handoff loss and speed feedback. DevOps privilege must still obey least privilege, separation of duties, MFA, monitoring, and controlled deployment.}
\noterow{DevSecOps}{Extends the integrated model so security/compliance are built into development and operations. Favor actionable, iterative improvement and evidence such as testing/red-team results; shift security earlier while retaining risk-based governance. There is no single standard implementation.}
\noterow{Modified Waterfall}{Adds limited flexibility to Waterfall: some next-phase work may begin before every current-phase task completes, and teams may step back a phase to correct assumptions and iterate.}
\noterow{Spiral}{An iterative, risk-centered model suited to large, complex, costly projects. Each cycle: (1) determine objectives/alternatives/risks; (2) evaluate alternatives and resolve risks; (3) develop/test; (4) plan the next phase. Repeated risk analysis improves adaptability but adds expense/overhead.}
\noterow{Prototype}{Build a simplified system, obtain stakeholder feedback, and iterate until needs are understood/met. Useful when requirements or a novel design cannot be reliably elicited in advance.}
\noterow{Cleanroom}{Focuses on defect prevention and high reliability rather than relying on post-development testing/remediation. Statistical measurement and rigorous design/implementation can support a defined reliability target, but at a cost to agility.}
\end{cornell}

\begin{exambox}{Method Selection Cue}
Waterfall optimizes predictability and phase control; Agile optimizes adaptation and frequent value delivery; Spiral explicitly repeats risk analysis; Prototype discovers uncertain requirements; Cleanroom emphasizes defect prevention. Security must be integrated whichever method governs the project.
\end{exambox}

\section{\color{navy}Maturity, Operations, Change, and Integrated Teams}
\source{pp. 561--571}
\begin{cornell}
\noterow{Why maturity models?}{Measure current capability, define a target state, identify gaps, and guide process improvement. More mature processes are generally more formal, repeatable, measured, and managed; maturity reduces some process risk but \textbf{does not guarantee project success}.}
\noterow{Common maturity elements}{\textbf{Domains} group practices; \textbf{levels} show increasing capability; \textbf{criteria} define what must be satisfied; \textbf{targets} express the desired state. Improvement should be justified by business benefit rather than an assumption that the highest level is always best.}
\noterow{CMM / CMMI}{CMM originated from DoD/SEI work on software capability. CMMI generalizes maturity across organizational process areas. CMMI levels: \textbf{1 Initial} (unpredictable/reactive), \textbf{2 Managed} (project-characterized), \textbf{3 Defined} (organization-wide/proactive), \textbf{4 Quantitatively Managed} (measured/controlled), \textbf{5 Optimizing} (continuous improvement).}
\noterow{OWASP SAMM}{The \textbf{Software Assurance Maturity Model} is prescriptive for building/improving a software security program and integrating security in the SDLC. Domains are Governance, Design, Implementation, Verification, and Operations; assess current state, choose target state, then improve people/process/knowledge/tools.}
\noterow{BSIMM}{The \textbf{Building Security-In Maturity Model} is empirically oriented: observed software-security initiatives show practices used in real organizations. It relies on a software security group (SSG) and organizes activities across Governance, Intelligence, Software Security Development Lifecycle, and Deployment.}
\noterow{CMMC}{DoD's Cybersecurity Maturity Model Certification addresses cybersecurity maturity for contractors handling controlled unclassified information (CUI). The book describes maturity progressing from basic performed practices to optimized, proactive capability and notes the role of independent assessment.}
\noterow{Operations and maintenance}{Production operation begins after verification/validation and deployment, but security work continues: monitoring, patching, vulnerability management, backups, incident response, access review, configuration management, and maintenance can trigger new SDLC iterations.}
\noterow{Change management}{Maintain known secure state while permitting necessary change. A formal CMB/CCB/CAB reviews requests, risk/impact, resources, scheduling, testing, approval, implementation, verification, documentation, and lessons learned; emergency changes use expedited control followed by documentation/review.}
\noterow{Integrated Product Team (IPT)}{A multidisciplinary team shares responsibility for delivering a product/process. Including users/customers, developers, security, operations, contractors, and other stakeholders brings relevant constraints/risks into decisions earlier and reduces waste/rework through faster feedback.}
\end{cornell}

\section{\color{navy}Security Controls in Development Ecosystems}
\source{pp. 572--590}
\begin{cornell}
\noterow{Programming-language risk}{Languages and paradigms influence memory management, type handling, error behavior, compilation/interpreting, privileges, and available safeguards. Object-oriented concepts include encapsulation, inheritance, polymorphism, and polyinstantiation; security depends on how objects/data and inherited behavior are controlled.}
\noterow{Markup and scripting}{HTML and CSS structure/present web content; JavaScript provides client-side interactivity; Python is a high-level interpreted language. Even simple/mobile code can be untrusted, so control sources, privileges, data access, validation, and execution environments.}
\noterow{Libraries / SDKs}{Reusable code accelerates development and can improve reliability, but a vulnerable or malicious library propagates risk to every dependent application. Approve trusted sources/versions, patch and inventory dependencies, and validate internally developed, vendor, or open-source libraries. SDKs bundle code/functions and development support for a platform.}
\noterow{Toolsets}{Editors, repositories/version control, debuggers/bug databases, compilers, and static/unit testing tools form the development toolchain. Govern approved versions/modules, configuration and patch state, access, changes, and developer training. A compromised tool can compromise many outputs.}
\noterow{IDE}{An \textbf{Integrated Development Environment} combines editing, debugging, build/compile, testing, and often deployment integrations. IDE policy/configuration can embed secure defaults and tests close to the developer; protect extensions, credentials, updates, and configuration.}
\noterow{Runtime / TCB}{A runtime includes the hardware, firmware, drivers, OS/interpreter, application, and services needed for execution. A \textbf{Trusted Computing Base (TCB)} is the hardware/software/firmware responsible for enforcing security and must be carefully designed, tested, isolated, and change-controlled.}
\noterow{Continuous Integration (CI)}{Continuously merge developer changes into a shared repository. Smaller/frequent integrations reduce large merge/debug events. Branching, pull requests, peer/automated reviews, and security tests protect the controlled baseline.}
\noterow{Continuous delivery vs. deployment}{\textbf{Delivery}: automatically moves tested code to a nonproduction stage or a production-ready state where manual action/verification may remain. \textbf{Deployment}: automatically releases successfully tested changes to production without manual intervention. Both are often abbreviated CD.}
\noterow{CI/CD pipeline}{The pipeline connects IDE, repository, build/package, testing, and deployment environments. Protect code and artifacts, secrets, service identities, definition files, integrations/APIs, logs, approvals, and environment boundaries; insert integrity, encryption, scanning, and review controls at appropriate stages.}
\noterow{IaC / continuous security}{\textbf{Infrastructure as Code (IaC)} definition files automate repeatable provisioning/configuration. Version/test/audit the definitions and embed automated AppSec checks such as unit, functional, regression, and vulnerability testing so security feedback follows each change.}
\noterow{SOAR}{\textbf{Security Orchestration, Automation, and Response} combines disparate security data/tools, executes playbooks/rules, and coordinates response. SIEM chiefly collects/correlates/alerts; SOAR additionally orchestrates and automates actions. Tune automation because false positives can cause denial of service.}
\noterow{Software Configuration Management}{SCM preserves integrity/visibility of source code, libraries, requirements, settings, and other \textbf{configuration items (CIs)}. Baselines represent known-good groups of CIs; changes occur in controlled/versioned branches, pass testing/QA/acceptance, then merge into the baseline.}
\noterow{Code repositories}{Central repositories support collaboration, branching, version history, access control, review, rollback, and automated testing. Treat source code and the repository as high-value assets: protect confidentiality, integrity, availability, access, backups, privileged changes, and audit history.}
\noterow{SAST}{\textbf{Static Application Security Testing} analyzes source/binaries without running the application. Fits IDE/check-in workflows and avoids production impact; drawbacks include language dependence and inability to observe many runtime/system-interaction flaws.}
\noterow{DAST}{\textbf{Dynamic Application Security Testing} exercises a running application through its interface/API, usually independent of source language. It can reveal runtime behavior, but findings may be farther from the developer and production testing can affect performance or discover flaws only after deployment.}
\noterow{IAST}{\textbf{Interactive Application Security Testing} combines source-aware and dynamic/runtime observations (and sometimes penetration-style techniques) to correlate findings. Tool/language support, expertise, cost, and performance are practical constraints.}
\noterow{RASP}{\textbf{Runtime Application Self-Protection} runs with the application to detect unexpected behavior and take corrective action such as blocking requests or virtual patching. It is protective rather than purely a test; false positives can disrupt service and compatibility requires care.}
\end{cornell}

\section{\color{navy}Assessing Software Security Effectiveness}
\source{pp. 590--598}
\begin{cornell}
\noterow{Auditing vs. logging}{\textbf{Logging} records events. \textbf{Auditing} systematically inspects activity/controls against expectations or standards. Logs become useful evidence only when protected and actually monitored/audited; capture what happened, actor/source, timestamp, and relevant metadata.}
\noterow{Accountability / nonrepudiation}{Accountability identifies the actor responsible for an event; nonrepudiation provides evidence strong enough that a unique actor cannot credibly deny the action. Not every log supports nonrepudiation; identity, integrity, time, and evidence strength matter.}
\noterow{Monitoring applications}{SIEM and other monitoring use logs for change verification, security-incident detection, performance/availability, cloud-service status, and investigations. Synchronize time, protect log integrity/confidentiality, define retention, and collect sufficient application/API/user context.}
\noterow{Development auditing}{Audit requirements artifacts (including classification/privacy analysis), design risk decisions and residual acceptance, implementation/code/testing evidence, deployment approvals/configuration, and operational changes. Adapt audit evidence collection to the methodology so short Agile iterations are not missed.}
\noterow{Risk-analysis focus}{Apply normal risk concepts across development and production. Prioritize sensitive/regulated data, evaluate new-technology uncertainty, reassess system changes, and control code risk. Threats can enter via tools, dependencies, repositories, people, build pipelines, environments, and acquired components.}
\noterow{SDLC risk sources}{Planning affects schedule/resources/security effort; requirements introduce attack/compliance needs; design selects architectures/providers/controls; development requires secure tools/coding/testing; O\&M carries live-use exposure and continuous patching, monitoring, and incident risks.}
\noterow{Mitigation decisions}{Integrate chosen controls into the SDLC and compare direct/indirect costs with risk reduction. Options include mitigation, transfer, compensating controls, or residual acceptance. Quantitative data are valuable where defensible; qualitative impacts help where outcomes are uncertain. Communicate assumptions and alternatives to management.}
\end{cornell}

\section{\color{navy}Security Impact of Acquired Software}
\source{pp. 599--604}
\begin{cornell}
\noterow{Acquired-software principle}{Buying/outsourcing reduces direct control and evidence over the developer's SDLC but can provide skills, speed, support, and capability the organization cannot economically build. Assess the source, obtain assurance, specify contractual controls, perform acceptance testing, and use compensating safeguards.}
\noterow{COTS}{\textbf{Commercial off-the-shelf} software provides established functionality/support but usually hides source code and developer practices and limits customer fixes. Maintain inventory, monitor vendor vulnerability/patch information, test externally where permitted, and rely on corrective/compensating controls when defect prevention is unavailable.}
\noterow{Open-source software (OSS)}{Source visibility permits independent review and community contribution, but visibility is not proof of review. Evaluate project/community health and security, validate source/repository integrity, assess support, inventory dependencies, patch promptly, and do not assume many reviewers guarantee flaw discovery.}
\noterow{Third-party custom software}{A specialized contractor can reduce skills/maturity risk but remains outside direct control. Put security, SDLC evidence, testing/acceptance, SLA, vulnerability/remediation, IP/source, and support obligations in the contract. \textbf{Code escrow} can release source under defined events such as supplier failure.}
\noterow{Managed services}{SaaS, PaaS, and IaaS shift different software responsibilities to a cloud service provider. The customer still determines whether the service is suitable for its data/use and retains responsibility for its data obligations. Direct assessment generally decreases as the provider manages more of the stack.}
\noterow{IaaS / PaaS / SaaS assessment}{\textbf{IaaS}: customer deploys more software and has more direct assessment/control. \textbf{PaaS}: provider manages more of the platform; assurance reports/SLAs compensate for lost control. \textbf{SaaS}: least direct application assessment; rely heavily on provider evidence, contracts, documented controls, and independent assurance.}
\noterow{Cloud assurance}{Third-party audit/assurance can reduce the need for every customer to directly inspect a CSP. The chapter points to SOC, CSA STAR, and ISO 27000-family assurance as examples; evaluate scope, period, exceptions, complementary customer controls, and fit to the organization's risk.}
\end{cornell}

\section{\color{navy}Secure Coding, Weaknesses, and Vulnerability Frameworks}
\source{pp. 604--612}
\begin{cornell}
\noterow{Standards vs. guidelines}{A \textbf{standard} is mandatory direction; a \textbf{guideline} is recommended direction. Secure-coding documentation converts policy objectives into developer behavior and can also form the baseline for code review, automated scanning, training, and audit.}
\noterow{Weakness vs. vulnerability}{A \textbf{weakness} is a flaw/condition that may be nonsecurity-related or not practically exploitable. A \textbf{vulnerability} is an exploitable security weakness in context. Compensating controls may make a source-level weakness nonexploitable, but that conclusion must be risk-assessed.}
\noterow{CWE}{\textbf{Common Weakness Enumeration} is a common language/catalog of software/hardware weakness types. CWE entries describe the weakness, relationships, introduction phase/platform, consequences/examples, observed CVEs, mitigations, and detection methods. It helps prevention, identification, tooling, and communication.}
\noterow{CWSS}{\textbf{Common Weakness Scoring System} prioritizes weaknesses. The book groups factors into \textbf{Base Finding}, \textbf{Attack Surface}, and \textbf{Environmental} groups. Treat the score as input to organization-specific analysis rather than an automatic remediation decision.}
\noterow{CVE}{\textbf{Common Vulnerabilities and Exposures} assigns standardized identifiers to disclosed vulnerabilities. An identifier includes registration year and numeric ID; records describe affected software/versions, risk, fixes/workarounds, and references. CVE identifies; it does not itself provide severity.}
\noterow{CVSS}{\textbf{Common Vulnerability Scoring System} scores severity from 0--10. Book categories: None 0; Low 0.1--3.9; Medium 4.0--6.9; High 7.0--8.9; Critical 9.0--10. Base metrics cover exploitability/impact; temporal metrics change with exploit/remediation confidence; environmental metrics tailor impact/exploitability to the organization.}
\noterow{CWE / CVE / CVSS}{CWE = \emph{type of weakness}; CVE = \emph{identifier for a disclosed vulnerability}; CVSS = \emph{severity/context scoring}. They complement one another and should be combined with asset criticality, exposure, threat intelligence, and compensating controls.}
\noterow{OWASP Top 10 (2017) 1--5}{As presented in this 2021 CBK: A1 Injection; A2 Broken Authentication; A3 Sensitive Data Exposure; A4 XML External Entities (XXE); A5 Broken Access Control. Use the list as developer/testing guidance, not as a CVE catalog or scoring system.}
\noterow{OWASP Top 10 (2017) 6--10}{A6 Security Misconfiguration; A7 Cross-Site Scripting (XSS); A8 Insecure Deserialization; A9 Using Components with Known Vulnerabilities; A10 Insufficient Logging \& Monitoring. Prevention requires secure design/configuration, validation/encoding, dependency governance, testing, and monitoring.}
\noterow{Software Composition Analysis}{\textbf{SCA} inventories software dependencies, identifies known vulnerabilities in those components, and supports update/remediation decisions. Integrate dependency checks into inventory and CI/CD; an automatic update can fix a flaw but also introduce availability/functionality risk, so test and govern it.}
\end{cornell}

\section{\color{navy}API Security}
\source{pp. 613--618}
\begin{cornell}
\noterow{API security model}{An API exposes standardized system functions/data to a consumer. REST commonly uses HTTP verbs/endpoints; SOAP and RPC are other forms. Treat exposed functions and all incoming API data as part of the attack surface and derive controls from the data/function risk.}
\noterow{Authentication}{Options include HTTP basic authentication (only with protected transport; Base64 is not encryption), API/key-based schemes, certificate-based authentication (including mutual authentication), and federated/SSO technologies such as SAML/OAuth. Authenticate the consumer, then authorize only necessary actions/data.}
\noterow{Authorization}{Provide granular control appropriate to user/system identities, roles, resources, functions, and data. Do not assume a system-to-system or developer-facing endpoint is trusted merely because humans do not normally interact with it.}
\noterow{Input validation / sanitization}{Treat external data as untrusted. Validate syntax and semantically valid values where practical; sanitize/escape content using technology-appropriate methods so data cannot become executable instructions. Escaping is particularly important against injection; semantic validation also protects data integrity.}
\noterow{Resource protection}{Threat-model the API just as rigorously as a UI. Limit exposed functionality, request/resource sizes and rates where relevant, permissions, and unintended information return. Test privileged/developer features because they can become attack paths.}
\noterow{Protected communications}{Use secure transport such as TLS according to sensitivity and authenticate endpoints. Protect confidentiality/integrity of data in transit and apply secure key/certificate lifecycle management.}
\noterow{Cryptography}{Prefer trusted/tested cryptographic modules to custom crypto. Encryption can provide confidentiality; hashes support integrity; HMAC can support authenticity/integrity; key ownership and signatures can support nonrepudiation/authenticity; key access is itself an access-control concern.}
\noterow{Logging / monitoring / alerting}{Log successful/failed authentication, input-validation outcomes, consumer/source, timestamp, operations, and security-relevant failures without exposing secrets. Review or ingest logs into monitoring/SIEM, correlate sources, and alert/escalate events that violate requirements.}
\noterow{API security testing}{Select functional/security tests based on API technology, data sensitivity, exposure, and functions. Use adequate coverage before penetration testing alone; include authorization, validation, error handling, transport, rate/resource protection, and logging behavior.}
\end{cornell}

\section{\color{navy}Secure Coding Practices and Software-Defined Security}
\source{pp. 618--624}
\begin{cornell}
\noterow{Secure-coding program}{Document, teach, and enforce organization-specific practices throughout development. Tailor recognized starting points such as OWASP and CMU SEI guidance to languages, architecture, risk, and technology; embed practices in IDE/tooling, code review, testing, and training.}
\noterow{OWASP practice areas}{The book's OWASP quick-reference categories cover input validation; output encoding; authentication/passwords; session management; access control; cryptography; error handling/logging; communication security; system configuration; database security; file management; memory management; and general coding practices.}
\noterow{Security culture}{Policies alone do not create secure software. Translate expectations into standards, procedures, checklists, education, and everyday decisions. Make security goals visible to developers and stakeholders and measure adherence/effectiveness.}
\noterow{Secure by design}{Address security in architecture/design before implementation. Assume attack will occur, identify threats/requirements early, and choose intrinsic controls rather than relying on expensive/error-prone retrofits.}
\noterow{Secure by default}{Ship conservative, protective configurations and layered controls as the starting state; require deliberate action to weaken protection or make resources public. Defaults matter because misconfiguration is common.}
\noterow{Shift left / build security in}{DevSecOps pushes security activities earlier in the lifecycle. SAMM/BSIMM can measure current capability and guide enterprise improvement. Early feedback reduces the cost and delay of discovering defects after deployment.}
\noterow{Software-Defined Security (SDS/SDSec)}{Virtualizes and software-controls security infrastructure, using definition/configuration and integrations to provision, monitor, manage, and reconfigure controls. It complements SDN/cloud/virtualized architectures where traditional perimeter or host agents may not fit.}
\noterow{SDS advantages}{The chapter highlights enhanced availability, infrastructure neutrality, micro-targeting/granularity, centralized management, standardization/integration/automation, and dynamic response. Definition files and orchestration make controls portable/repeatable but also make their integrity and privilege high-value security concerns.}
\noterow{Dynamic response}{Automated reconfiguration can rapidly isolate hosts, alter virtual network paths/policies, or deploy compensating protection after an event. Treat ``self-healing'' claims cautiously: automated actions require rules, testing/tuning, monitoring, governance, and human oversight.}
\end{cornell}

\section{\color{navy}Critical Comparisons}
\begin{center}\small
\begin{tabularx}{\textwidth}{>{\bfseries\color{navy}\raggedright\arraybackslash}p{0.23\textwidth}>{\raggedright\arraybackslash}X>{\raggedright\arraybackslash}X}
\toprule Concept & First idea & Distinguishing idea\\\midrule
Waterfall / Agile & Sequential, phase-gated & Iterative, change-responsive\\
DevOps / DevSecOps & Integrate development + operations & Explicitly integrate security/compliance\\
CI / delivery / deployment & Merge changes frequently & Production-ready with possible gate / automatic production release\\
SAST / DAST & Inspect nonrunning code/artifacts & Exercise a running application/interface\\
IAST / RASP & Correlate source + runtime testing & Runtime protection/action\\
CWE / CVE / CVSS & Weakness type & Vulnerability ID / severity score\\
COTS / OSS & Vendor binary/source mostly hidden & Source visible/community-development model\\
Standard / guideline & Mandatory direction & Recommended direction\\
\bottomrule
\end{tabularx}
\end{center}

\section{\color{navy}Key Terms and Acronyms}
\begin{summarybox}{Rapid-Review Glossary}
\textbf{BSIMM} -- Building Security-In Maturity Model \quad
\textbf{CI/CD} -- Continuous Integration / Continuous Delivery or Deployment \quad
\textbf{CMMI} -- Capability Maturity Model Integration \quad
\textbf{COTS} -- Commercial Off-the-Shelf \quad
\textbf{CVE/CVSS/CWE/CWSS} -- Vulnerability/weakness identification and scoring systems \quad
\textbf{DAST/IAST/SAST/RASP} -- Application security testing/protection approaches \quad
\textbf{IaC} -- Infrastructure as Code \quad
\textbf{IDE} -- Integrated Development Environment \quad
\textbf{OSS/FOSS} -- (Free and) Open-Source Software \quad
\textbf{SAMM} -- Software Assurance Maturity Model \quad
\textbf{SCA} -- Software Composition Analysis \quad
\textbf{SCM} -- Software Configuration Management \quad
\textbf{SDLC} -- Software Development Life Cycle \quad
\textbf{SDS/SDSec} -- Software-Defined Security \quad
\textbf{SOAR} -- Security Orchestration, Automation, and Response \quad
\textbf{TDD} -- Test-Driven Development.
\end{summarybox}

\section{\color{navy}High-Yield Review}
\begin{itemize}
\item Security belongs throughout the SDLC and broader SLC; a final security test cannot compensate for missing requirements and insecure design.
\item Agile is principles; XP, TDD, and Scrum are methods/practices that implement iterative ideas. DevSecOps integrates security into DevOps delivery.
\item CMMI levels run Initial $\rightarrow$ Managed $\rightarrow$ Defined $\rightarrow$ Quantitatively Managed $\rightarrow$ Optimizing; maturity does not guarantee success.
\item Continuous \emph{delivery} may retain a manual production gate; continuous \emph{deployment} automatically releases successful changes.
\item SAST sees code without execution; DAST sees runtime externally; IAST combines perspectives; RASP acts at runtime.
\item CWE describes weakness types, CVE identifies disclosed vulnerabilities, and CVSS scores severity; organization context still determines priority.
\item Acquired/cloud software trades direct control for capability. Assurance, contracts, acceptance testing, and compensating controls manage the lost visibility.
\item API inputs are untrusted; authenticate, authorize, validate/sanitize, encrypt, log, monitor, and test exposed functions.
\end{itemize}

\section{\color{navy}Active-Recall Questions}
\begin{enumerate}
\item How does the SDLC differ from the broader SLC, and where should security be integrated?
\item Compare Waterfall, Agile, Spiral, Prototype, and Cleanroom on change/risk/defect handling.
\item Name the three key Scrum roles and explain backlog versus sprint.
\item What does DevSecOps add to DevOps, and which classic security principles still apply to privileged delivery teams?
\item List the five CMMI maturity levels in order.
\item Distinguish CI, continuous delivery, and continuous deployment.
\item Compare SAST, DAST, IAST, and RASP.
\item What is the difference between logging and auditing, and between accountability and nonrepudiation?
\item What security trade-offs distinguish COTS, OSS, third-party custom software, and SaaS/PaaS/IaaS?
\item Distinguish CWE, CVE, CVSS, and SCA.
\item Which controls are fundamental to securing an API?
\item What advantages does software-defined security provide, and what new concentration of risk does automation create?
\end{enumerate}

\subsection*{\color{teal}Answer Key}
{\small
\begin{enumerate}[nosep]
\item SDLC focuses on creating/evolving software; SLC includes operation/maintenance and disposal. Integrate security requirements, design, implementation, testing, deployment, operation, change, and retirement activities rather than one gate.
\item Waterfall is sequential; Agile iterates and accepts change; Spiral iterates risk analysis; Prototype iterates stakeholder feedback to discover needs; Cleanroom emphasizes rigorous defect prevention/reliability.
\item Scrum master, product owner, and development team. The backlog is prioritized remaining product work; a sprint is a time-boxed committed subset/delivery cycle.
\item It explicitly integrates security/compliance and actionable security feedback. Least privilege, separation of duties, MFA, monitoring, testing/approval, and controlled changes still apply.
\item Initial, Managed, Defined, Quantitatively Managed, Optimizing.
\item CI merges changes frequently; delivery keeps code deployable and can retain a manual production step; deployment automatically releases successful code to production.
\item SAST analyzes nonrunning source/artifacts; DAST executes the app externally; IAST combines source/runtime context; RASP monitors and blocks/mitigates while the app runs.
\item Logging records events; auditing systematically reviews them/controls. Accountability attributes an event; nonrepudiation provides evidence strong enough to prevent a unique actor from denying it.
\item COTS offers vendor support but low source/control visibility; OSS offers source visibility/community but no automatic assurance; custom third-party work relies on contract/acceptance/escrow; cloud services shift increasing stack responsibility to the provider from IaaS toward SaaS.
\item CWE catalogs weakness types; CVE identifies specific disclosed vulnerabilities; CVSS scores vulnerability severity/context; SCA inventories/analyzes third-party/open-source dependencies for vulnerability risk.
\item Risk-based authentication, granular authorization, input validation/sanitization, least exposure/resource protection, protected communications/crypto, secure key management, logging/monitoring/alerting, and appropriate functional/security testing.
\item SDS enables portable/consistent policy, automation, micro-targeting, centralized management, infrastructure neutrality, availability, and rapid response. Definition files, orchestration, service privilege, and automation can also replicate mistakes or compromise at scale.
\end{enumerate}
}

\section{\color{navy}Cornell Summary}
\begin{summarybox}{Domain 8 Summary}
Secure software comes from a governed lifecycle, not a last-minute scan. Choose a development method suited to the work, integrate security stakeholders and tests early, mature the process deliberately, and protect the entire toolchain from IDE through repository, build, CI/CD, and runtime. Measure effectiveness with logs, audits, testing, and risk analysis. When software or platforms are acquired, replace lost direct control with evidence, contracts, testing, and compensating controls. Use secure coding standards, dependency analysis, CWE/CVE/CVSS context, strong API controls, and software-defined automation to reduce risk while keeping human governance over high-impact decisions.
\end{summarybox}

\section{\color{navy}Final Domain Checklist}
{\small
\begin{itemize}[label=\(\square\),nosep]
\item I can compare the major SDLC methodologies and integrate security into each lifecycle phase.
\item I can explain CMMI, SAMM, BSIMM, and the purpose/limits of maturity models.
\item I can secure programming ecosystems, toolchains, repositories, CI/CD, SCM, and runtime components.
\item I can distinguish SAST, DAST, IAST, and RASP and choose complementary testing approaches.
\item I can use auditing/logging and risk analysis to assess software-security effectiveness.
\item I can assess COTS, OSS, custom third-party, and cloud-managed software under shared responsibility.
\item I can distinguish CWE, CWSS, CVE, CVSS, OWASP Top 10, and software composition analysis.
\item I can apply API security, secure-coding standards, and software-defined security concepts.
\end{itemize}
}
\vfill{\footnotesize\color{teal}Prepared as a study companion to Deane \& Kraus, \textit{The Official (ISC)\textsuperscript{2} CISSP CBK Reference}, 6th ed. Content is paraphrased and synthesized for study.}
\end{document}
